\chapter{IPv4 Addressing and Subnetting}\label{ch:ipv4}
\bp{1.3}

\begin{outcomes}
By the end of this chapter you will be able to:
\begin{tight}
  \item \textbf{Convert} between dotted decimal, binary and prefix length fast
        enough to do it in an exam without a calculator.
  \item \textbf{Compute} the network address, broadcast address, usable range and
        host count of any subnet, in under thirty seconds.
  \item \textbf{Subnet} a given block to meet a stated requirement, and apply
        variable-length masks without overlapping.
  \item \textbf{Diagnose} the four addressing faults that account for almost every
        broken host: wrong mask, wrong gateway, address outside the subnet, and
        duplicate address.
\end{tight}
\end{outcomes}

\begin{prereq}
\chref{models} and \chref{ethernet}. You need to be comfortable that a switch
works on MAC addresses and a router works on IP addresses.
\end{prereq}

\begin{keyterms}
dotted decimal $\bullet$ prefix length $\bullet$ subnet mask $\bullet$ network
address $\bullet$ broadcast address $\bullet$ usable range $\bullet$ block size
$\bullet$ magic number $\bullet$ RFC 1918 $\bullet$ private address $\bullet$
public address $\bullet$ VLSM $\bullet$ default gateway $\bullet$ APIPA
\end{keyterms}

\begin{examnote}[This topic is worth more time than any other in Part I]
Topic \tp{1.3} says \emph{troubleshoot} IPv4 addressing, assignment and
subnetting, ``public and private''. You will be given a host that cannot reach
something and expected to spot that its mask is wrong. Every second you spend
computing a subnet boundary is a second not spent reasoning, so the arithmetic
has to be automatic. Learn the method in §\ref{sec:magic} until it is boring.
\end{examnote}

\section{Counting in binary, once}

Everything in this chapter is a consequence of one fact: an IPv4 address is
\textbf{32 bits}, and we write it as four decimal numbers only for our own
convenience. The shortcut in Section~\ref{sec:magic} lets you subnet without ever
writing a bit down --- but it will be a magic trick you can forget rather than a
method you can rebuild unless you look at the bits properly once. This is that
once.

\begin{definitionbox}[Bits, bytes and octets]
A \term{bit} is a single 0 or 1. Eight bits make a \term{byte}, which in
networking is almost always called an \term{octet} --- the words mean the same
thing here.

Each position in an octet is worth twice the one to its right:
\end{definitionbox}

% Nine columns: the budget is about 12.8 cm, not the 15.78 of two.
\begin{longtable}{@{}L{2.9cm}L{1.2cm}L{1.2cm}L{1.2cm}L{1.2cm}L{1.2cm}L{1.2cm}L{1.2cm}L{1.2cm}@{}}
\toprule
\rowcolor{primary!15}
\textbf{Position value} &
\textbf{128} & \textbf{64} & \textbf{32} & \textbf{16} &
\textbf{8} & \textbf{4} & \textbf{2} & \textbf{1} \\
\midrule
\endhead
\cmd{192} is & \cmd{1} & \cmd{1} & \cmd{0} & \cmd{0} & \cmd{0} & \cmd{0} & \cmd{0} & \cmd{0} \\
\rowcolor{lightbg}
\cmd{255} is & \cmd{1} & \cmd{1} & \cmd{1} & \cmd{1} & \cmd{1} & \cmd{1} & \cmd{1} & \cmd{1} \\
\cmd{77} is & \cmd{0} & \cmd{1} & \cmd{0} & \cmd{0} & \cmd{1} & \cmd{1} & \cmd{0} & \cmd{1} \\
\bottomrule
\end{longtable}

\begin{conceptbox}[Both directions, in one line each]
\textbf{Binary to decimal:} add the position values wherever there is a 1.
$64+8+4+1 = 77$.

\textbf{Decimal to binary:} walk left to right, and at each position ask ``does
this fit in what is left?''. For 77: 128 no (0), 64 yes (1, leaving 13), 32 no
(0), 16 no (0), 8 yes (1, leaving 5), 4 yes (1, leaving 1), 2 no (0), 1 yes (1,
leaving 0). That gives \cmd{01001101}.

Two facts fall straight out and are worth holding on to. \textbf{255 is all eight
bits set} --- which is why it is the largest value an octet can hold, and why a
mask octet of 255 means ``all of this octet is network''. And \textbf{0 is all
eight clear}, meaning ``none of it is''.
\end{conceptbox}

\begin{worked}[Where the block size comes from]
Take \cmd{192.168.10.77} with mask \cmd{255.255.255.192}, and write only the
last octet of each in bits:

\begin{center}
\begin{tabular}{@{}ll@{}}
address \cmd{77} & \cmd{01\,001101} \\
mask \cmd{192}   & \cmd{11\,000000} \\
\end{tabular}
\end{center}

The mask has \textbf{two} 1 bits in this octet, so two bits of it are network and
the remaining \textbf{six} are host. Set those six host bits to zero and you have
the network address; set them all to one and you have the broadcast address:

\begin{center}
\begin{tabular}{@{}lll@{}}
network   & \cmd{01\,000000} & $= 64$ \\
broadcast & \cmd{01\,111111} & $= 127$ \\
\end{tabular}
\end{center}

So the subnet is \cmd{192.168.10.64} to \cmd{192.168.10.127} --- and it is
\textbf{64 addresses} wide, because six free host bits give $2^6 = 64$
combinations.

Now notice: $256 - 192 = 64$ as well. That is not a coincidence and it is not a
trick. Subtracting the mask value from 256 is simply a fast way of asking ``how
many combinations do the remaining host bits have?'' \textbf{That is the magic
number}, and Section~\ref{sec:magic} does the whole job with it.
\end{worked}

\begin{conceptbox}[Powers of two, which you should know on sight]
Host bits decide subnet size, so these eleven numbers cover almost every question
you will be asked.

\medskip
\begin{tabular}{@{}lllllllllll@{}}
$2^1$ & $2^2$ & $2^3$ & $2^4$ & $2^5$ & $2^6$ & $2^7$ & $2^8$ & $2^9$ & $2^{10}$ & $2^{11}$ \\
2 & 4 & 8 & 16 & 32 & 64 & 128 & 256 & 512 & 1024 & 2048 \\
\end{tabular}

\medskip
\textbf{Usable} hosts is always two fewer than the total, because the all-zeros
address names the network itself and the all-ones address is the broadcast.
Neither can be given to a host. So six host bits give 64 addresses and
\textbf{62} usable ones.
\end{conceptbox}

\section{The address and the mask}

An IPv4 address is 32 bits, written as four decimal bytes. It carries two pieces
of information at once, and the mask is what separates them.

\begin{definitionbox}[Network part and host part]
The \term{subnet mask} marks, with 1 bits, which leading bits of the address
identify the \term{network}. The remaining 0 bits identify the \term{host} within
it. Written as a count of 1 bits, the mask is a \term{prefix length}:
\cmd{255.255.255.0} and \cmd{/24} are the same statement.

Two addresses are on the same subnet if and only if their network parts are
identical. That single sentence is the source of most addressing faults: a host
with the wrong mask computes a different network part, decides its neighbour is
remote, and sends the traffic to a gateway that may not want it.
\end{definitionbox}

\begin{definitionbox}[Default gateway]
The address a host sends traffic to when the destination is \textbf{not} on its
own subnet. Almost always a router or a layer 3 switch interface on that subnet.

Every host makes exactly one decision for every packet it sends, and it makes it
with the mask:

\begin{tight}
  \item \textbf{Same network part as mine?} Then the destination is local. ARP for
        it directly and send the frame straight there. The gateway is not
        involved at all.
  \item \textbf{Different?} Then it is remote. ARP for the \emph{gateway},
        and hand the frame to it --- with the destination's IP address still
        inside, but the gateway's MAC address on the outside.
\end{tight}

Three consequences that account for a great many faults, and that
\chref{clienttshoot} will make you test for directly:

\begin{tight}
  \item \textbf{A host with no gateway can reach its own subnet and nothing
        else.} Everything local works, so the user reports ``the internet is
        down'' rather than ``my machine is misconfigured''.
  \item \textbf{The gateway must be inside the host's own subnet.} A host will
        not ARP for something it believes is remote, so a gateway outside the
        subnet is unreachable by definition --- which is the fault in this
        chapter's Break \& Fix.
  \item \textbf{Get the mask wrong and the host makes the local-or-remote
        decision wrongly}, for some destinations and not others. That is why a
        wrong mask presents as ``some things work'', which is far harder to read
        than ``nothing works''.
\end{tight}
\end{definitionbox}

\begin{longtable}{@{}C{1.6cm}L{3.6cm}C{2.0cm}C{2.4cm}L{4.6cm}@{}}
\toprule
\rowcolor{primary!15}
\textbf{Prefix} & \textbf{Mask} & \textbf{Block} & \textbf{Usable} & \textbf{Typical use} \\
\midrule
\endhead
/24 & 255.255.255.0   & 256 & 254 & A standard user VLAN \\
\rowcolor{lightbg}
/25 & 255.255.255.128 & 128 & 126 & Half a /24 \\
/26 & 255.255.255.192 & 64  & 62  & A small office \\
\rowcolor{lightbg}
/27 & 255.255.255.224 & 32  & 30  & A rack or a server segment \\
/28 & 255.255.255.240 & 16  & 14  & A tiny segment \\
\rowcolor{lightbg}
/29 & 255.255.255.248 & 8   & 6   & Four routers and room to grow \\
/30 & 255.255.255.252 & 4   & 2   & A point-to-point link between routers \\
\rowcolor{lightbg}
/31 & 255.255.255.254 & 2   & 2   & A point-to-point link, no waste (RFC 3021) \\
/32 & 255.255.255.255 & 1   & 1   & A host route or a loopback \\
\bottomrule
\end{longtable}

Usable is block size minus two, because the first address in a block is the
network address and the last is the broadcast address. The exception is /31,
where both are usable because a point-to-point link has no need for a broadcast.

\begin{definitionbox}[Private addresses --- RFC 1918]
Three ranges are reserved for internal use and are never routed on the internet.
The exam expects you to recognise them instantly.
\begin{tight}
  \item \cmd{10.0.0.0/8} --- one very large block
  \item \cmd{172.16.0.0/12} --- that is \cmd{172.16.0.0} to \cmd{172.31.255.255},
        and the upper bound catches people out
  \item \cmd{192.168.0.0/16}
\end{tight}
Also worth knowing: \cmd{169.254.0.0/16} is \term{APIPA}, which a host assigns
itself when DHCP fails. A host with a 169.254 address has not been given one ---
see \chref{dhcp}.
\end{definitionbox}

\section{The method: block size, or ``magic number''}\label{sec:magic}

You have now seen where this comes from, so it is a shortcut rather than a spell.
In the exam you will not have time to write out binary --- and after this section
you will not need to.

\begin{conceptbox}[Four steps, every time]
\begin{tightnum}
  \item \textbf{Find the interesting octet} --- the one where the mask is neither
        255 nor 0.
  \item \textbf{Block size} $= 256 -$ (the mask value in that octet). This is the
        magic number.
  \item \textbf{Count in blocks} from zero in that octet until you pass the
        address. The multiple you last passed is the \textbf{network address}.
  \item \textbf{Broadcast} $=$ next network address $-1$. Usable range is
        everything between the two, exclusive.
\end{tightnum}
\end{conceptbox}

\begin{worked}[Where does 172.16.83.190/26 live?]
\begin{tightnum}
  \item Mask for /26 is \cmd{255.255.255.192}. The interesting octet is the
        fourth.
  \item Block size $= 256 - 192 = 64$.
  \item Count in 64s in the fourth octet: 0, 64, 128, \textbf{192}. Our host is
        190, which sits after 128 and before 192. Network address is
        \cmd{172.16.83.128}.
  \item Next network is \cmd{172.16.83.192}, so broadcast is
        \cmd{172.16.83.191}. Usable range is \cmd{172.16.83.129} to
        \cmd{172.16.83.190}.
\end{tightnum}
So \cmd{.190} is the \emph{last usable address} in its subnet --- a perfectly
valid host address, and a favourite exam distractor because it looks like a
broadcast.
\end{worked}

\begin{worked}[The same method in the third octet: 10.44.150.7/21]
\begin{tightnum}
  \item /21 is \cmd{255.255.248.0}. Interesting octet is the third.
  \item Block size $= 256 - 248 = 8$.
  \item Count in 8s in the third octet: 0, 8, \ldots, 144, \textbf{152}. Our 150
        sits after 144. Network is \cmd{10.44.144.0}.
  \item Next network is \cmd{10.44.152.0}, so broadcast is
        \cmd{10.44.151.255}. Usable is \cmd{10.44.144.1} to
        \cmd{10.44.151.254} --- $2046$ hosts.
\end{tightnum}
Note that the fourth octet plays no part in finding the boundary. Only the
interesting octet does.
\end{worked}

\section{Subnetting a block to a requirement}

\begin{worked}[Split 192.168.20.0/24 for four sites of 50, 25, 10 and 2 hosts]
Work \emph{largest first}, or you will strand addresses.

\begin{tightnum}
  \item 50 hosts needs 62 usable $\rightarrow$ /26, block 64.
        \cmd{192.168.20.0/26} --- usable \cmd{.1}--\cmd{.62}.
  \item 25 hosts needs 30 usable $\rightarrow$ /27, block 32. Next free address
        is \cmd{.64}: \cmd{192.168.20.64/27} --- usable \cmd{.65}--\cmd{.94}.
  \item 10 hosts needs 14 usable $\rightarrow$ /28, block 16. Next free is
        \cmd{.96}: \cmd{192.168.20.96/28} --- usable \cmd{.97}--\cmd{.110}.
  \item 2 hosts (a router link) $\rightarrow$ /30, block 4. Next free is
        \cmd{.112}: \cmd{192.168.20.112/30} --- usable \cmd{.113}--\cmd{.114}.
\end{tightnum}
Everything from \cmd{.116} upwards is still free. This is \term{VLSM}: different
masks inside one block, sized to the requirement.
\end{worked}

\begin{alertbox}[The two rules VLSM must obey]
\begin{tight}
  \item \textbf{A subnet must start on a multiple of its own block size.} A /27
        can start at \cmd{.64} or \cmd{.96}, never at \cmd{.70}.
  \item \textbf{Subnets must not overlap.} If you allocate smallest-first you
        will eventually place a large subnet where it does not fit, and have to
        start again. Largest first avoids this.
\end{tight}
\end{alertbox}

\section{Configuring and verifying an address}

\begin{config}{an interface address and a default gateway}
! On a router interface -- the mask is written in full.
interface GigabitEthernet0/0
 description User VLAN 10 gateway
 ip address 192.168.20.1 255.255.255.192
 no shutdown
!
! On a layer 2 switch, the management address lives on an SVI,
! and the switch needs a default gateway because it does not route.
interface Vlan10
 ip address 192.168.20.2 255.255.255.192
 no shutdown
!
ip default-gateway 192.168.20.1
\end{config}

\begin{verify}{R1}{show ip interface brief}
Interface              IP-Address      OK? Method Status                Protocol
GigabitEthernet0/0     192.168.20.1    YES manual up                    up
GigabitEthernet0/1     unassigned      YES unset  administratively down down
Vlan1                  unassigned      YES unset  down                  down
\end{verify}

\begin{conceptbox}[Reading the two status columns]
The pair of columns is the single most useful diagnostic on a router, and the
combinations mean different things.
\begin{tight}
  \item \textbf{up / up} --- working.
  \item \textbf{up / down} --- the cable and electrical link are fine, but the
        protocol is not. Usually an encapsulation or keepalive mismatch, or the
        far end is not configured.
  \item \textbf{down / down} --- a physical problem. Cable, transceiver, or the
        other end is shut. Go back to \chref{physical}.
  \item \textbf{administratively down / down} --- somebody typed
        \cmd{shutdown}. The fix is \cmd{no shutdown}.
\end{tight}
\end{conceptbox}

\begin{pitfall}
\begin{tight}
  \item \textbf{Forgetting \cmd{no shutdown}.} Router interfaces are shut by
        default. Switch ports are not. This asymmetry catches everyone once.
  \item \textbf{Assuming \cmd{172.16.0.0/12} stops at 172.16.} It runs to
        \cmd{172.31.255.255}. \cmd{172.32.0.0} is public.
  \item \textbf{Giving the gateway an address outside the subnet.} A host will
        accept the configuration and then fail silently for everything off-net.
  \item \textbf{Using the network or broadcast address on a host.} \cmd{.0} and
        \cmd{.255} are only network and broadcast in a /24. In a /26 the
        boundaries move, and \cmd{192.168.20.64} is a network address even though
        it does not end in zero.
\end{tight}
\end{pitfall}

\begin{breakfix}[``I can reach my own floor but nothing else''.]
A host is configured \cmd{192.168.20.100} mask \cmd{255.255.255.192}, gateway
\cmd{192.168.20.1}. It can ping \cmd{192.168.20.90} but not its gateway and
nothing beyond.

\textbf{Diagnosis.} Work the mask. /26 gives block size 64, so the boundaries are
\cmd{.0}, \cmd{.64}, \cmd{.128}, \cmd{.192}. The host at \cmd{.100} is in the
\cmd{192.168.20.64/26} subnet, usable \cmd{.65}--\cmd{.126}. Its gateway at
\cmd{.1} is in the \emph{first} subnet, \cmd{192.168.20.0/26}. The host and its
gateway are on different subnets, so the host cannot ARP for it and has no way
off its own segment. It can still reach \cmd{.90} because that address is in the
same /26.

\textbf{Fix.} Either move the host into the subnet its gateway serves, or --- far
more likely what was intended --- correct the gateway to \cmd{192.168.20.65}, the
first usable address of the subnet the host is actually in. Then confirm with a
ping to the gateway before testing anything remote.
\end{breakfix}

\begin{lab}{Subnet, configure, then break it}{Packet Tracer}
\textbf{Topology.} One router, two switches, four PCs --- two per VLAN.

\begin{tasks}
  \item Subnet \cmd{172.16.5.0/24} into two subnets of at least 100 hosts each.
        Write down, before touching a device: network address, first and last
        usable address, and broadcast address for both.
  \item Configure the router with a gateway address in each subnet and the PCs
        with addresses from your plan. Verify with \cmd{show ip interface brief}
        and a ping within each subnet.
  \item Prove routing works between the two subnets by pinging across.
  \item Now change one PC's mask from your planned value to
        \cmd{255.255.255.0}. Predict what will break \emph{before} you test, then
        test it. Explain the result in terms of the network part.
  \item Change it back, then set one PC's gateway to an address in the other
        subnet. Predict and test again.
  \item \textbf{Extension.} Re-plan the same /24 using VLSM for four segments of
        100, 50, 20 and 2 hosts, and configure it without overlaps.
\end{tasks}
\end{lab}

\begin{checkpoint}
\begin{tightnum}
  \item For \cmd{10.19.200.37/27}: give the network address, broadcast address,
        and first and last usable address.
  \item A host is \cmd{192.168.1.130/25} and its gateway is
        \cmd{192.168.1.126}. Will it work? Show your reasoning.
  \item How many usable addresses does a /30 provide, and why is a /31 sometimes
        used instead on a router-to-router link?
\end{tightnum}
\end{checkpoint}

\begin{chaptersummary}
\begin{tight}
  \item Address plus mask equals two facts: which network, and which host.
        Devices are on the same subnet only if their network parts match.
  \item Block size $=256-$ mask octet. Count in blocks. That method answers every
        subnetting question you will be asked.
  \item Usable $=$ block $-2$, except /31 and /32.
  \item Private ranges are \cmd{10/8}, \cmd{172.16/12} (to 172.31) and
        \cmd{192.168/16}. A \cmd{169.254} address means DHCP failed.
  \item Plan VLSM largest first, and start every subnet on a multiple of its own
        block size.
  \item Four faults cover most cases: wrong mask, wrong gateway, address outside
        the subnet, duplicate address.
\end{tight}
\end{chaptersummary}

\begin{reviewq}
  \item \textbf{(MCQ)} A host is addressed \cmd{192.168.90.200/28}. What is the
        broadcast address of its subnet?
        (a)~192.168.90.207 (b)~192.168.90.255 (c)~192.168.90.191 (d)~192.168.90.215
  \item \textbf{(MCQ)} Which is a valid host address in \cmd{10.0.8.0/22}?
        (a)~10.0.8.0 (b)~10.0.11.255 (c)~10.0.11.254 (d)~10.0.12.1
  \item \textbf{(Short)} Explain why a host with the correct address but a mask
        that is too short may still reach some destinations and not others.
  \item \textbf{(Short)} A host has address \cmd{169.254.14.9}. What has
        happened, and what would you check first?
  \item \textbf{(Applied)} You are given \cmd{192.168.40.0/24} and asked for six
        subnets: 60, 60, 28, 12, 12 and 2 hosts. Produce a VLSM plan with no
        overlaps, showing network address, mask and usable range for each, and
        state how much of the block remains free.
\end{reviewq}
